% !TEX program = xelatex
% ═══════════════════════════════════════════════════════════════════
%  \en{ZiqaKernel} — بحث علمي أكاديمي شامل
%  Architectural & Systems Analysis: Academic Research Paper
%  Language: Arabic (RTL) — XeLaTeX + polyglossia + FreeSerif
% ═══════════════════════════════════════════════════════════════════
\documentclass[12pt,a4paper]{article}

% ── اللغة والخطوط ────────────────────────────────────────────────
\usepackage{fontspec}

% ── الرياضيات ────────────────────────────────────────────────────
\usepackage{amsmath,amssymb,amsthm,mathtools}

% ── الجداول ─────────────────────────────────────────────────────
\usepackage{booktabs,array,multirow,tabularx,longtable,colortbl}
\newcolumntype{C}[1]{>{\centering\arraybackslash}p{#1}}
\newcolumntype{R}[1]{>{\raggedleft\arraybackslash}p{#1}}
\newcolumntype{L}[1]{p{#1}}

% ── الرسوميات ───────────────────────────────────────────────────
\usepackage{graphicx,float,wrapfig,subcaption}

% ── TikZ وpgfplots ───────────────────────────────────────────────
\usepackage{tikz}
\usepackage{tikz-3dplot}
\usepackage{pgfplots}
\usepackage{pgf-pie}
\usetikzlibrary{
  positioning, shapes, shapes.geometric, arrows.meta,
  decorations.pathreplacing, decorations.markings,
  calc, 3d, perspective,
  shadows.blur, patterns, fit,
  matrix, backgrounds, mindmap, trees,
  fadings, shadings
}
\pgfplotsset{compat=1.18}
\usepgfplotslibrary{colormaps}

% ── الألوان ─────────────────────────────────────────────────────
\usepackage{xcolor}
\definecolor{accent}{HTML}{2C3E50} % Dark Slate Blue
\definecolor{accent2}{HTML}{2980B9} % Professional Blue
\definecolor{accent3}{HTML}{C0392B} % Professional Red
\definecolor{accent4}{HTML}{8E44AD} % Muted Purple
\definecolor{accentg}{HTML}{D35400} % Muted Orange
\definecolor{codebg}{HTML}{F8F9FA} % Very light gray
\definecolor{titlebg}{HTML}{2C3E50} % Dark Slate Blue
\definecolor{rowalt}{HTML}{F2F2F2}
\definecolor{darkgray}{HTML}{34495E}
\definecolor{lightblue}{HTML}{E8F6F3}
\definecolor{deepblue}{HTML}{2C3E50}
\definecolor{teal}{HTML}{16A085}
\definecolor{purple}{HTML}{8E44AD}
\definecolor{orange}{HTML}{D35400}

% ── تحسينات مرئية ───────────────────────────────────────────────
\usepackage{tcolorbox}
\tcbuselibrary{skins, breakable, theorems, most}

% صناديق ملونة للتعريفات والنظريات
\newtcolorbox{definitionbox}[1][]{
  enhanced, colback=white, colframe=accent, arc=0pt,
  boxrule=0.5pt, left=8pt, right=8pt, top=6pt, bottom=6pt,
  title={#1}, fonttitle=\bfseries\normalsize\color{white},
  attach boxed title to top left={yshift=-2mm,xshift=2mm},
  boxed title style={colback=accent, arc=0pt, boxrule=0pt}
}

\newtcolorbox{resultbox}[1][]{
  enhanced, colback=white, colframe=accent2, arc=0pt,
  boxrule=0.5pt, left=8pt, right=8pt, top=6pt, bottom=6pt,
  title={#1}, fonttitle=\bfseries\normalsize\color{white},
  attach boxed title to top left={yshift=-2mm,xshift=2mm},
  boxed title style={colback=accent2, arc=0pt, boxrule=0pt}
}

\newtcolorbox{warningbox}[1][]{
  enhanced, colback=white, colframe=accentg, arc=0pt,
  boxrule=0.5pt, left=8pt, right=8pt, top=6pt, bottom=6pt,
  title={#1}, fonttitle=\bfseries\normalsize\color{white},
  attach boxed title to top left={yshift=-2mm,xshift=2mm},
  boxed title style={colback=accentg, arc=0pt, boxrule=0pt}
}

% ── الكود ───────────────────────────────────────────────────────
\usepackage{listings}
\lstdefinestyle{rustcode}{
  language={},
  backgroundcolor=\color{codebg},
  basicstyle=\ttfamily\footnotesize,
  breaklines=true, frame=single,
  rulecolor=\color{accent!40},
  numbers=left, numberstyle=\tiny\color{gray},
  tabsize=4, xleftmargin=2em, xrightmargin=1em,
  escapeinside={(*@}{@*)},
  morekeywords={fn,let,mut,struct,enum,impl,match,if,else,for,
    while,return,Ok,Err,Some,None,pub,use,mod,self,super,
    crate,alloc,vec,\en{Result},Option,true,false,unsafe,async,await},
  keywordstyle=\color{accent}\bfseries,
  commentstyle=\color{accent2}\itshape\arabicfont,
  stringstyle=\color{accent3},
  showstringspaces=false,
}
\lstdefinestyle{ccode}{
  language=C,
  backgroundcolor=\color{codebg},
  basicstyle=\ttfamily\footnotesize,
  breaklines=true, frame=single,
  rulecolor=\color{accent!40},
  numbers=left, numberstyle=\tiny\color{gray},
  tabsize=4, xleftmargin=2em, xrightmargin=1em,
  escapeinside={(*@}{@*)},
  keywordstyle=\color{accent}\bfseries,
  commentstyle=\color{accent2}\itshape\arabicfont,
  stringstyle=\color{accent3},
}
\lstdefinestyle{pseudo}{
  language={},
  backgroundcolor=\color{deepblue!8},
  basicstyle=\ttfamily\footnotesize,
  breaklines=true, frame=tb,
  rulecolor=\color{deepblue!50},
  tabsize=2, xleftmargin=2em,
  mathescape=true,
}

% ── تخطيط الصفحة والمسافات ────────────────────────────────────────────────
\usepackage{geometry}
\geometry{a4paper, top=2.5cm, bottom=2.5cm, left=2.5cm, right=2.5cm}
\usepackage{setspace}
\setstretch{1.35} % تباعد الأسطر الأكاديمي
\setlength{\parskip}{0.4em plus 0.1em minus 0.1em}
\setlength{\parindent}{1.5em}

\tolerance=5000
\emergencystretch=3em
\hbadness=5000

\usepackage{fancyhdr}
\usepackage{titlesec}
\usepackage{caption}
\usepackage{enumitem}
\setlist{itemsep=0.3em, topsep=0.4em, parsep=0.1em, partopsep=0pt}
% ── الروابط ─────────────────────────────────────────────────────

\usepackage[numbers,sort&compress]{natbib}
\usepackage[colorlinks=true, linkcolor=accent,
  citecolor=accent2, urlcolor=teal, unicode=true]{hyperref}
\setlength{\headheight}{32pt}
\pagestyle{fancy}
\fancyhf{}
\fancyhead[L]{\small\color{darkgray}%
  بحث علمي: تحليل معماري لنواة \textsc{ZiqaKernel} --- \textit{يونيو 2026}}
\fancyhead[R]{\small\color{darkgray}\thepage}
\fancyhead[L]{\small\color{darkgray}%
  بحث علمي: تحليل معماري لنواة \en{ZiqaKernel} --- \textit{يونيو 2026}}
\fancyhead[R]{\small\color{darkgray}\thepage}
\fancyfoot[C]{\small\color{gray!70}\rule{0.4\textwidth}{0.4pt}}
\renewcommand{\headrulewidth}{0.4pt}
\renewcommand{\headrule}{\hbox to\headwidth{%
  \color{gray!30}\leaders\hrule height\headrulewidth\hfill}}
\titleformat{\section}
  {\Large\bfseries\color{titlebg}}
  {\thesection\hspace{8pt}}
  {0em}{}[\vspace{4pt}\color{titlebg}\rule{\textwidth}{0.8pt}]

\titleformat{\subsection}
  {\large\bfseries\color{accent}}
  {\thesubsection\hspace{8pt}}
  {0em}{}

\titleformat{\subsubsection}
  {\normalsize\bfseries\color{darkgray}}
  {\textcolor{accent2}{\small$\blacksquare$}\hspace{6pt}\thesubsubsection\hspace{6pt}}
  {0em}{}
% ── البيئات الرياضية ─────────────────────────────────────────────
\theoremstyle{definition}
\newtheorem{definition}{تعريف}[section]
\newtheorem{theorem}{نظرية}[section]
\newtheorem{lemma}{مبرهنة}[section]
\newtheorem{corollary}{نتيجة}[section]
\newtheorem{remark}{ملاحظة}[section]

% ── أوامر مساعدة ────────────────────────────────────────────────
\newcommand{\tier}[1]{\textbf{\textcolor{accent2}{#1-Tier}}}
\newcommand{\keyword}[1]{\textbf{\textsc{#1}}}
\renewcommand{\arraystretch}{1.3}
\newcommand{\codeinline}[1]{\texttt{\small #1}}
\newcommand{\highlight}[1]{\colorbox{lightblue}{\small #1}}

% ══════════════════════════════════════════════════════════════════
% ── اللغة ───────────────────────────────────────────────────
\usepackage{polyglossia}
\setmainlanguage[numerals=maghrib]{arabic}
\setotherlanguage{english}
\newfontfamily\arabicfont[Script=Arabic,Scale=1.15,WordSpace=1.5,
  UprightFont=Amiri-Regular,BoldFont=Amiri-Bold,ItalicFont=Amiri-Italic]{Amiri}
\newfontfamily\arabicfontsf[Script=Arabic,Scale=1.1,
  UprightFont=Amiri-Regular,BoldFont=Amiri-Bold]{Amiri}
\newfontfamily\arabicfonttt[Script=Arabic,Scale=1.0,
  UprightFont=Amiri-Regular]{Amiri}
\setmainfont[
  UprightFont=Amiri-Regular,BoldFont=Amiri-Bold,ItalicFont=Amiri-Italic]{Amiri}
\setmonofont[Script=Latin,Language=English]{Liberation Mono}

\newcommand{\ar}[1]{\textarabic{#1}}
\newcommand{\Ar}[1]{\textarabic{#1}}
\newcommand{\en}[1]{\textenglish{#1}}

\begin{document}

% ── صفحة العنوان ─────────────────────────────────────────────────
\begin{titlepage}
\begin{tikzpicture}[remember picture, overlay]
  % خلفية متدرجة
  \fill[left color=titlebg, right color=deepblue!80]
    (current page.north west) rectangle
    ([yshift=-9cm]current page.north east);
  % زخرفة هندسية
  \foreach \i in {0,0.4,...,3.6}{
    \draw[white, opacity=0.05, line width=0.8pt]
      ([xshift=\i cm, yshift=-0.5cm]current page.north west)
      -- ([xshift=\i cm+18cm, yshift=-8.5cm]current page.north west);
  }
  \foreach \i in {0,0.6,...,8}{
    \draw[white, opacity=0.05, line width=0.8pt]
      ([yshift=-\i cm]current page.north west)
      -- ([yshift=-\i cm+0.5cm]current page.north east);
  }
  % دائرة زخرفية
  \draw[white, opacity=0.12, line width=2pt] 
    ([xshift=14cm,yshift=-4cm]current page.north west) circle(3.5cm);
  \draw[white, opacity=0.06, line width=4pt] 
    ([xshift=14cm,yshift=-4cm]current page.north west) circle(4.8cm);
  % شريط سفلي
  \fill[accent2] ([yshift=-9cm]current page.north west)
    rectangle ([yshift=-9.5cm]current page.north east);
  \fill[accentg, opacity=0.7] ([yshift=-9.5cm]current page.north west)
    rectangle ([yshift=-9.7cm]current page.north east);
\end{tikzpicture}

\vspace*{1.5cm}
\begin{center}
{\fontsize{11}{14}\selectfont\color{lightblue}\bfseries%
  المستودع العلمي الشامل --- يونيو 2026}

\vspace{0.6cm}

{\fontsize{32}{38}\selectfont\bfseries\color{white}%
  تحليل معماري ونظامي}

\vspace{0.3cm}

{\fontsize{26}{32}\selectfont\bfseries\color{lightblue}%
  لنواة \textsc{ZiqaKernel}}

\vspace{0.5cm}
{\color{white!80}\rule{0.6\textwidth}{0.6pt}}
\vspace{0.4cm}

{\fontsize{13}{18}\selectfont\color{white!90}%
  \textit{\en{Microkernel} Design · \en{Capability-Based} Security}\\[4pt]
  \textit{\en{eBPF Obsidian-Tier} Verification · Real-Time Bare Metal Applications}}

\vspace{0.5cm}
\end{center}

\vspace{0.8cm}

% بطاقات الإحصاءات
\begin{center}
\begin{tikzpicture}
  \def\cardw{2.8cm}
  \def\cardh{2.0cm}
  % بطاقة 1
  \node[draw=white!40, fill=white!15, rounded corners=8pt, blur shadow={shadow blur steps=5, shadow opacity=40},
        minimum width=\cardw, minimum height=\cardh, text=white] at (0,0) {
    \begin{minipage}{2.5cm}\centering
      {\fontsize{22}{26}\bfseries 28,633}\\[2pt]
      {\footnotesize\textarabic{سطور كود Rust}}
    \end{minipage}};
  % بطاقة 2
  \node[draw=white!40, fill=white!15, rounded corners=8pt, blur shadow={shadow blur steps=5, shadow opacity=40},
        minimum width=\cardw, minimum height=\cardh, text=white] at (3.2,0) {
    \begin{minipage}{2.5cm}\centering
      {\fontsize{22}{26}\bfseries 2,015}\\[2pt]
      {\footnotesize\textarabic{عقد المعرفة}}
    \end{minipage}};
  % بطاقة 3
  \node[draw=white!40, fill=white!15, rounded corners=8pt, blur shadow={shadow blur steps=5, shadow opacity=40},
        minimum width=\cardw, minimum height=\cardh, text=white] at (6.4,0) {
    \begin{minipage}{2.5cm}\centering
      {\fontsize{22}{26}\bfseries 3,843}\\[2pt]
      {\footnotesize\textarabic{علاقة معمارية}}
    \end{minipage}};
  % بطاقة 4
  \node[draw=white!40, fill=white!15, rounded corners=8pt, blur shadow={shadow blur steps=5, shadow opacity=40},
        minimum width=\cardw, minimum height=\cardh, text=white] at (9.6,0) {
    \begin{minipage}{2.5cm}\centering
      {\fontsize{22}{26}\bfseries 131}\\[2pt]
      {\footnotesize\textarabic{مجتمعاً معمارياً}}
    \end{minipage}};
  % بطاقة 5
  \node[draw=white!40, fill=white!15, rounded corners=8pt, blur shadow={shadow blur steps=5, shadow opacity=40},
        minimum width=\cardw, minimum height=\cardh, text=white] at (12.8,0) {
    \begin{minipage}{2.5cm}\centering
      {\fontsize{22}{26}\bfseries 101}\\[2pt]
      {\footnotesize\textarabic{ملف مصدر}}
    \end{minipage}};
\end{tikzpicture}
\end{center}

\vspace{1.2cm}

\begin{center}
\begin{tabular}{C{3.2cm}C{3.2cm}C{3.2cm}C{3.2cm}}
\rowcolor{white!20}
\bfseries\color{lightblue}\textarabic{نموذج الأمان} &
\bfseries\color{lightblue}\textarabic{المعمارية} &
\bfseries\color{lightblue}\textarabic{اللغة} &
\bfseries\color{lightblue}\textarabic{مستوى التطوير} \\[4pt]
\color{white} Capability A+ &
\color{white} x86\_64 &
\color{white} Rust + Zig &
\color{white} مُصلَّب
\end{tabular}

\vspace{1.5cm}

{\small\color{white!70}%
  الباحث: \textbf{\en{S}.~Shamsi} \quad|\quad
  المنهجية: \textsc{graphify} Knowledge Graph Pipeline\quad|\quad
  \LaTeX{} + \textsc{XeLaTeX}}
\end{center}

\vfill
\end{titlepage}

% ══════════════════════════════════════════════════════════════════
% الملخص
% ══════════════════════════════════════════════════════════════════
\newpage
\section*{الملخص}

\begin{tcolorbox}[enhanced, colback=accent!6, colframe=accent,
  arc=4pt, boxrule=1.2pt, left=10pt, right=10pt, top=8pt, bottom=8pt]

\textbf{\large ملخص عربي:}\quad
تقدم هذه الورقة البحثية تحليلاً معمارياً ونظامياً شاملاً لنواة
\textsc{ZiqaKernel}، وهي نظام تشغيل تجريبي من نوع
\en{microkernel} مكتوب بلغة \textsc{Rust} مع دمج مختار لمسارات الأداء الحرجة بلغة
\textsc{Zig}، مصممة للعمل على المعدن العاري لمعمارية \texttt{x86\_64}.

\medskip

تم تحليل ما يزيد على 28,633 سطراً من الكود الموزع على 101 ملف مصدر، طُوِّرت خلال 5 أيام من البحث المكثف بمعدل 60 التزاماً في مستودع git. تقدم النواة نموذج أمان فريداً قائماً على القدرات (\en{capability-based}) مع إلغاء فوري ومتكرر للتفويض عبر شجرة إلغاء، وجدولة \en{MLFQ} \en{A-Tier} مع 4 طوابير وأولوية ديناميكية، ونظام \en{eBPF Obsidian-Tier} مع مفسّر متكامل ومدقق يستخدم التحليل التجريدي (\en{abstract interpretation}).

\medskip

استُخرجت خريطة معرفية كاملة للمشروع باستخدام خط أنابيب \textsc{graphify}، منتجةً 2,015 عقدة و3,843 علاقة موزعة على 131 مجتمعاً معمارياً. كشفت النتائج عن نضج تقني في جميع المكونات الأساسية للنواة مع تحقيق جميع الأولويات P0 و P1.

\bigskip\hrule\bigskip

\textbf{\large Abstract:}\quad
This paper presents a comprehensive architectural and systems analysis of the
\textsc{ZiqaKernel}, an experimental \en{microkernel} operating system written in
\textsc{Rust} with selective \textsc{Zig} integration for performance-critical
paths, targeting \texttt{x86\_64} bare metal. Over 28{,}633 lines of code
spanning 101 source files across 5 days of intensive development were analyzed.
The kernel introduces a unique \en{capability-based} security model with immediate
revocation, an \en{A-Tier} \en{MLFQ} \en{scheduler}, and an \en{Obsidian-Tier} eBPF subsystem with
\en{abstract interpretation} verification. A complete knowledge graph of 2{,}015
nodes and 3{,}843 edges across 131 architectural communities was extracted via
the \textsc{graphify} pipeline. All P0 and P1 priorities are achieved.
\end{tcolorbox}

\vspace{0.5cm}

\noindent\textbf{الكلمات المفتاحية:}
نواة تجريبية، \en{Microkernel}، Rust، Zig، Capability Security،
\en{MLFQ} \en{Scheduler}، eBPF, \en{Abstract Interpretation}، x86\_64، Knowledge Graph

\noindent\textbf{Keywords:}
\en{Microkernel}, Rust, Zig, \en{Capability-Based} Security, \en{MLFQ} Scheduling,
eBPF Verification, \en{Abstract Interpretation}, Bare Metal, Knowledge Graph

\newpage
\tableofcontents
\newpage

% ══════════════════════════════════════════════════════════════════
\section{مقدمة}
% ══════════════════════════════════════════════════════════════════

تمثل نواة \textsc{ZiqaKernel} مختبراً بحثياً تجريبياً لاستكشاف حدود تصميم أنظمة التشغيل الحديثة. تُكتب النواة أساساً بلغة \textsc{Rust}---وهي لغة تضمن سلامة الذاكرة في وقت الترجمة---مع دمج انتقائي لمسارات حرجة بلغة \textsc{Zig} لتحقيق أقصى أداء. المشروع بالغ التعقيد: 28,633 سطر كود، 101 ملف مصدر، 5 أيام تطوير مكثف.

\subsection{الدوافع البحثية}

تنبثق أهمية هذا البحث من ثلاثة إشكاليات جوهرية تعاني منها أنظمة التشغيل الحديثة:

\begin{enumerate}[label=\arabic*., leftmargin=2em]
\item \textbf{أزمة سلامة الذاكرة:} ما يزال نحو \en{70\%} من الثغرات الأمنية الحرجة ناجماً عن أخطاء إدارة الذاكرة في أنظمة مكتوبة بلغة C. لغة Rust تحل هذه الإشكالية بضمانات وقت الترجمة.
\item \textbf{جمود نموذج الأمان:} نماذج الأمان التقليدية القائمة على المستخدمين والمجموعات (DAC/MAC) تعاني من ضعف بنيوي. نموذج القدرات (\en{Capability-Based}) يوفر تحكماً أدق وإلغاءً فورياً.
\item \textbf{التوتر بين الأداء والمرونة:} الجداول الثابتة لا تتكيف مع أعباء العمل المتغيرة. \en{MLFQ} يحقق توازناً أمثل بين الاستجابية والإنتاجية.
\end{enumerate}

\subsection{الإسهامات الرئيسية للبحث}

\begin{definitionbox}[الإسهامات الرئيسية]
\begin{enumerate}[label=\textbf{C\arabic*.}, leftmargin=3em]
\item تصميم وتنفيذ نظام قدرات (Capability System) مع إلغاء فوري ومتكرر عبر شجرة إلغاء متعددة المستويات---صُنِّف \en{A+ Tier}.
\item جدولة \en{MLFQ} مع 4 طوابير وأولوية ديناميكية ومنع التجويع---صُنِّف \en{A-Tier}.
\item آلة eBPF افتراضية كاملة مع مدقق يستخدم التحليل التجريدي---صُنِّف \en{Obsidian-Tier}.
\item طبقة \en{ABI} قابلة للتوسيع بـ 129 استدعاء نظام (99 متوافقة مع \en{Linux} \en{POSIX}).
\item تشغيل تطبيقات تفاعلية على المعدن العاري: \en{Shell} (40+ أمر)، \en{Editor}، \en{Tetris}، \en{DOOM}.
\item دمج \en{Rust/Zig FFI} عبر \en{C-ABI} لمشغل \en{blitter} بأداء \en{ReleaseFast}.
\item استخراج خريطة معرفة كاملة للمشروع بـ 2,015 عقدة و 3,843 علاقة عبر \textsc{graphify}.
\end{enumerate}
\end{definitionbox}

\subsection{تنظيم الورقة}

تُنظَّم بقية هذه الورقة كالتالي: القسم \ref{sec:method} يشرح منهجية التحليل. القسم \ref{sec:arch} يعرض المعمارية الكاملة. القسم \ref{sec:scheduler} يتناول الجدولة وإدارة الذاكرة. القسم \ref{sec:ebpf} يُحلل نظام eBPF. القسم \ref{sec:perf} يقيس الأداء. القسم \ref{sec:compare} يقارن مع أنوية أخرى. القسم \ref{sec:graph} يعرض خريطة المعرفة. القسم \ref{sec:conc} يختتم بالاستنتاجات.

\newpage

% ══════════════════════════════════════════════════════════════════
\section{منهجية التحليل}\label{sec:method}
% ══════════════════════════════════════════════════════════════════

تعتمد هذه الدراسة على منهجية تحليل ثلاثية المراحل مُنفَّذة عبر خط أنابيب \textsc{graphify}، تجمع بين التحليل الهيكلي الحتمي والاستخراج الدلالي والتجميع الطوبولوجي.

\subsection{مراحل خط أنابيب التحليل}

\begin{center}
\begin{tikzpicture}[
  node distance=1.2cm and 2.5cm,
  stage/.style={draw=accent, fill=accent!10, rounded corners=8pt,
    minimum width=3.5cm, minimum height=1.2cm,
    text width=3.2cm, align=center, font=\small\bfseries},
  arrow/.style={->, thick, accent, shorten >=2pt, shorten <=2pt},
  label/.style={font=\tiny, text=darkgray, align=center, text width=2.5cm}
]
  \node[stage] (s1) {الاستخراج الهيكلي\\[2pt]\footnotesize AST Analysis};
  \node[stage, right=of s1] (s2) {الاستخراج الدلالي\\[2pt]\footnotesize LLM Extraction};
  \node[stage, right=of s2] (s3) {بناء الرسم البياني\\[2pt]\footnotesize Graph Construction};
  \node[stage, right=of s3] (s4) {التحليل والتقرير\\[2pt]\footnotesize Louvain + Report};

  \draw[arrow] (s1) -- (s2);
  \draw[arrow] (s2) -- (s3);
  \draw[arrow] (s3) -- (s4);

  \node[label, below=0.3cm of s1] {Rust/Zig\\ محلل AST\\ (حتمي، مجاني)};
  \node[label, below=0.3cm of s2] {نماذج LLM\\ عمليات موازية\\ 20-25 ملف/وكيل};
  \node[label, below=0.3cm of s3] {دمج البيانات\\ رسم بياني واحد\\ 2015 عقدة};
  \node[label, below=0.3cm of s4] {خوارزمية Louvain\\ 131 مجتمع\\ God Nodes};
\end{tikzpicture}
\end{center}

\vspace{0.5cm}

\subsubsection{المرحلة الأولى: الاستخراج الهيكلي (\en{AST})}

تحليل شجري حتمي لكود المصدر يستخرج:
\begin{itemize}[leftmargin=2em]
\item تعريفات الدوال والهياكل والوحدات والأنواع
\item علاقات الاستيراد والاستدعاء
\item تعليقات الوثائق وتعليقات التوثيق (\en{doc comments})
\item توقيعات الأنواع والقيود
\end{itemize}

\subsubsection{المرحلة الثانية: الاستخراج الدلالي}

نماذج \en{LLM} تُعالج ملفات الوثائق والأوراق البحثية والصور. يُوزَّع الملفات على وكلاء متوازيين (20-25 ملفاً/وكيل) ثم تُدمج النتائج. لا تتطلب هذه المرحلة تعديل الكود.

\subsubsection{المرحلة الثالثة: البناء والتجميع}

يُدمج ناتج المرحلتين في رسم بياني واحد، ثم تُطبَّق خوارزمية لوفان (\en{Louvain}) لاكتشاف المجتمعات المعمارية. تُحلَّل العقد المركزية (\en{God Nodes}) وتُولَّد التقارير.

\subsection{عقد الله (\en{God Nodes}) --- أكثر المكونات ترابطاً}

جدول العقد المركزية الخمسة عشر الأولى في الرسم البياني المعرفي:

\vspace{0.3cm}
\begin{center}
\begin{tabular}{C{0.5cm}L{4cm}R{2cm}R{3cm}C{3cm}}
\toprule
\rowcolor{titlebg!90}
\color{white}\textbf{\#} &
\color{white}\textbf{الرمز} &
\color{white}\textbf{الدرجة} &
\color{white}\textbf{نسبة المركزية} &
\color{white}\textbf{الفئة} \\
\midrule
\rowcolor{rowalt} 1 & \en{Shell} & 76 & \en{12.1\%} & واجهة المستخدم \\
2 & \en{Result} & 74 & \en{11.8\%} & معالجة الأخطاء \\
\rowcolor{rowalt} 3 & \en{AbiError} & 74 & \en{11.8\%} & طبقة \en{ABI} \\
4 & \en{SyscallContext} & 74 & \en{11.8\%} & استدعاءات النظام \\
\rowcolor{rowalt} 5 & \en{String} & 54 & \en{8.6\%} & بنية بيانات \\
6 & \en{Editor} & 52 & \en{8.3\%} & واجهة المستخدم \\
\rowcolor{rowalt} 7 & \en{Scheduler} & 36 & \en{5.7\%} & جدولة العمليات \\
8 & run() & 36 & \en{5.7\%} & نقطة دخول \\
\rowcolor{rowalt} 9 & \en{Screen} & 26 & \en{4.1\%} & مدير العرض \\
10 & \en{ZiqaFs} & 23 & \en{3.7\%} & نظام الملفات \\
\rowcolor{rowalt} 11 & \en{BlockDevice} & 22 & \en{3.5\%} & سائق التخزين \\
12 & \en{CapabilitySpace} & 21 & \en{3.4\%} & نظام الأمان \\
\rowcolor{rowalt} 13 & \en{VirtioNet} & 19 & \en{3.0\%} & سائق الشبكة \\
14 & \en{Compositor} & 18 & \en{2.9\%} & مدير النوافذ \\
\rowcolor{rowalt} 15 & \en{PidAllocator} & 17 & \en{2.7\%} & إدارة العمليات \\
\bottomrule
\end{tabular}
\end{center}

\newpage

% ══════════════════════════════════════════════════════════════════
\section{المعمارية الشاملة}\label{sec:arch}
% ══════════════════════════════════════════════════════════════════

\subsection{نظرة عامة على الطبقات الثلاث}

يُبنى \textsc{ZiqaKernel} على هيكل هرمي ثلاثي الطبقات: مساحة المستخدم (\en{Ring 3})، نواة النظام (\en{Ring 0})، وطبقة العتاد. كل طبقة تتصل بالأخرى عبر واجهات محكمة الضبط.

\vspace{0.5cm}

% رسم معماري ثلاثي الأبعاد بواسطة tikz-3dplot
\begin{center}
\tdplotsetmaincoords{70}{25}
\begin{tikzpicture}[tdplot_main_coords, scale=0.9]

% ─── طبقة العتاد (Layer 0) ───
\begin{scope}[canvas is xy plane at z=0]
  \fill[orange!20, draw=orange!70, line width=1pt, rounded corners=3pt]
    (-6,-2.5) rectangle (6,2.5);
  \node[font=\bfseries\small, text=orange!80!black] at (0,2.0)
    {العتاد (Hardware)};
  % مكونات العتاد
  \foreach \x/\lbl in {-4.5/CPU,-2.5/RAM,-0.5/NVMe,1.5/VirtIO,3.5/VGA,5/ACPI}{
    \fill[orange!40, draw=orange!60, rounded corners=2pt]
      (\x-0.8,-1.6) rectangle (\x+0.8,-0.2);
    \node[font=\tiny, text=orange!90!black] at (\x,-0.9) {\lbl};
  }
\end{scope}

% ─── نواة النظام (Layer 1) ───
\begin{scope}[canvas is xy plane at z=3]
  \fill[accent!20, draw=accent!70, line width=1.2pt, rounded corners=3pt]
    (-6,-2.5) rectangle (6,2.5);
  \node[font=\bfseries\small, text=accent!90!black] at (0,2.0)
    {نواة النظام (Ring 0)};
  \foreach \x/\lbl in {-4.5/Microkernel,-2.5/{MLFQ},-0.5/VFS,
                        1.5/{eBPF VM},3.5/Capability,5/IPC}{
    \fill[accent!35, draw=accent!50, rounded corners=2pt]
      (\x-0.9,-1.6) rectangle (\x+0.9,-0.2);
    \node[font=\tiny, text=accent!90] at (\x,-0.9) {\lbl};
  }
\end{scope}

% ─── مساحة المستخدم (Layer 2) ───
\begin{scope}[canvas is xy plane at z=6]
  \fill[accent2!20, draw=accent2!70, line width=1.2pt, rounded corners=3pt]
    (-6,-2.5) rectangle (6,2.5);
  \node[font=\bfseries\small, text=accent2!90!black] at (0,2.0)
    {مساحة المستخدم (Ring 3)};
  \foreach \x/\lbl in {-4/{Shell (40+)}, -1.5/{Editor},
                        1/{Tetris}, 3.5/{DOOM}, 5.5/{libposix}}{
    \fill[accent2!35, draw=accent2!50, rounded corners=2pt]
      (\x-1.0,-1.6) rectangle (\x+1.0,-0.2);
    \node[font=\tiny, text=accent2!80!black] at (\x,-0.9) {\lbl};
  }
\end{scope}

% ─── الأسهم الرأسية ───
\foreach \x in {-5, 0, 5}{
  \draw[->, thick, accent!60!black, opacity=0.8]
    (\x, 0, 0.2) -- (\x, 0, 2.8);
  \draw[->, thick, accent2!60!black, opacity=0.8]
    (\x, 0, 3.2) -- (\x, 0, 5.8);
}

% ─── تسميات الجانب ───
\node[font=\scriptsize, text=accent2!80, rotate=90, anchor=south]
  at (-7.5, 0, 6) {تطبيقات};
\node[font=\scriptsize, text=accent!80, rotate=90, anchor=south]
  at (-7.5, 0, 3) {نواة};
\node[font=\scriptsize, text=orange!80, rotate=90, anchor=south]
  at (-7.5, 0, 0) {عتاد};

% ─── الأسطح الجانبية (3D effect) ───
\begin{scope}[opacity=0.15]
  \fill[accent2!50]
    (-6,-2.5,6) -- (-6,-2.5,5.8) -- (6,-2.5,5.8) -- (6,-2.5,6) -- cycle;
  \fill[accent!50]
    (-6,-2.5,3) -- (-6,-2.5,2.8) -- (6,-2.5,2.8) -- (6,-2.5,3) -- cycle;
\end{scope}

\end{tikzpicture}
\captionof{figure}{الرسم المعماري ثلاثي الأبعاد لطبقات \textsc{ZiqaKernel} --- من العتاد إلى مساحة المستخدم}
\end{center}

\subsection{وحدات هيكل الكود}

\begin{center}
\begin{tabular}{L{3.5cm}L{7.5cm}R{1.8cm}}
\toprule
\rowcolor{titlebg!85}
\bfseries\color{white}\textarabic{الوحدة} &
\bfseries\color{white}\textarabic{المحتوى} &
\bfseries\color{white}\textarabic{السطور} \\
\midrule
\rowcolor{rowalt}
\codeinline{src/arch/x86\_64/} &
\codeinline{gdt, idt, paging, interrupts, per\_cpu, switch, kvm} & 2,103 \\
\codeinline{src/memory/} &
\codeinline{frame, heap, paging, heapstats, page\_table} & 1,037 \\
\rowcolor{rowalt}
\codeinline{src/process/} &
\codeinline{\en{scheduler}, signal, vma, switch, pid} & 1,137 \\
\codeinline{src/ebpf/} &
\codeinline{verifier, vm, map, attach, mod} & 2,583 \\
\rowcolor{rowalt}
\codeinline{src/fs/} &
\codeinline{\en{ziqafs}, ramfs, vfs, page\_cache, fat32, file\_ops} & 2,837 \\
\codeinline{src/drivers/} &
\codeinline{virtio*, block, uart, keyboard, rtl8139, pci} & 3,590 \\
\rowcolor{rowalt}
\codeinline{src/\en{abi}/} &
\codeinline{syscall, \en{linux}, \en{wasm}, libposix} & 4,470 \\
\codeinline{src/userspace/} &
\codeinline{\en{shell}, \en{editor}, \en{tetris}, \en{doom}, drm\_driver} & 10,906 \\
\bottomrule
\end{tabular}
\end{center}

\subsection{توزيع الكود البياني}

\begin{center}
\begin{tikzpicture}
\begin{axis}[
  x dir=reverse,
  ybar,
  width=\textwidth, height=7cm,
  symbolic x coords={\en{ABI},Drivers,FS,eBPF,Sched.,Memory,Arch,User},
  xtick=data,
  xticklabel style={font=\small, rotate=30, anchor=east},
  ylabel={سطور الكود},
  ylabel style={font=\small},
  ymin=0, ymax=12000,
  nodes near coords,
  nodes near coords align={vertical},
  every node near coord/.append style={font=\footnotesize, rotate=90,
    anchor=west, yshift=3pt},
  bar width=14pt,
  enlarge x limits=0.08,
  ymajorgrids=true,
  grid style={dashed, gray!30},
  axis line style={gray!60},
  xtick style={color=gray!60},
]
\addplot[fill=accent!85, draw=accent, fill opacity=0.8, draw opacity=1] coordinates {
  (\en{ABI},4470) (Drivers,3590) (FS,2837)
  (eBPF,2583) (Sched.,1137) (Memory,1037)
  (Arch,2103) (User,10906)
};
\end{axis}
\end{tikzpicture}
\captionof{figure}{توزيع سطور الكود عبر الوحدات الرئيسية للنواة}
\end{center}

\vspace{0.5cm}

\begin{center}
\begin{tikzpicture}
\pie[
  text=pin,
  radius=3.0,
  style={drop shadow},
  color={accent!80, accent2!70, orange!70, purple!65,
         teal!70, accent3!65, accentg!70, accent!40}
]{
  15.6/\small \en{ABI} (4470),
  12.5/\small Drivers (3590),
  9.9/\small FS (2837),
  9.0/\small eBPF (2583),
  7.3/\small Arch+Sched+Mem (4277),
  38.0/\small User-Space (10906)
}
\end{tikzpicture}
\captionof{figure}{النسب المئوية لتوزيع الكود حسب الوحدة}
\end{center}

\newpage

% ══════════════════════════════════════════════════════════════════
\section{الجدولة وإدارة الذاكرة}\label{sec:scheduler}
% ══════════════════════════════════════════════════════════════════

\subsection{نظام القدرات (Capability System) --- تصنيف A+}

نظام القدرات هو المكوِّن الأمني الأساسي في \textsc{ZiqaKernel}. يتجاوز النماذج التقليدية ليقدم إلغاءً فورياً متكرراً عبر شجرة تفويض.

\begin{definitionbox}[نظام القدرات]
القدرة (Capability) هي رمز مميت مرتبط بعملية بعينها، يمنحها حق الوصول إلى مورد محدد بصلاحيات محددة. يمكن تفويض القدرة لعملية أخرى مع إمكانية الإلغاء الفوري في أي وقت، مما يُقطع الوصول على كل المستويات الفرعية بشكل آني.
\end{definitionbox}

\vspace{0.4cm}
رسم بياني لشجرة تفويض القدرات:

\begin{center}
\begin{tikzpicture}[
  every node/.style={draw, rounded corners=4pt, font=\small,
    minimum height=0.6cm, minimum width=1.8cm},
  root/.style={fill=accent, text=white, draw=accent!80},
  child1/.style={fill=accent2!70, draw=accent2},
  child2/.style={fill=accentg!60, draw=accentg},
  revoked/.style={fill=accent3!30, draw=accent3, dashed},
  active/.style={fill=accent2!30, draw=accent2},
  edge from parent/.style={draw, ->, thick}
]
\node[root] (root) {الجذر (نواة)}
  child[level distance=1.5cm, sibling distance=5cm] {
    node[child1] (p1) {Process A}
    child[level distance=1.3cm, sibling distance=2.5cm] {
      node[active] (c1) {Process B}
    }
    child[level distance=1.3cm] {
      node[revoked] (c2) {مُلغاة ×}
    }
  }
  child[level distance=1.5cm, sibling distance=5cm] {
    node[child2] (p2) {Process C}
    child[level distance=1.3cm, sibling distance=2.5cm] {
      node[active] (c3) {Process D}
    }
    child[level distance=1.3cm] {
      node[active] (c4) {Process E}
    }
  };
\node[right=0.5cm of c2, font=\footnotesize, text=accent3, draw=none]
  {إلغاء فوري};
\draw[->, red, dashed] (p1.east) to[bend right=20] (c2.north);
\end{tikzpicture}
\captionof{figure}{شجرة تفويض القدرات مع إلغاء فوري متكرر}
\end{center}

\vspace{0.3cm}

\begin{center}
\begin{tabular}{L{3cm}L{8.5cm}}
\toprule
\rowcolor{titlebg!85}
\bfseries\color{white}\textarabic{الخاصية} & \bfseries\color{white}\textarabic{التفاصيل} \\
\rowcolor{rowalt} النوع & \en{Capability-based} مع إلغاء فوري متكرر \\
التفويض & شجرة إلغاء متعددة المستويات (Revocation Tree) \\
\rowcolor{rowalt} القدرات & ملفات، أجهزة، قنوات IPC، شبكة \\
الاستدعاءات &
  \codeinline{cap\_close, cap\_revoke, dev\_port\_in/out, dev\_map, dev\_irq\_wait} \\
\rowcolor{rowalt} التصنيف & \textbf{\color{accent2} A+ Tier} \\
الضمانات & لا يمكن لأي عملية الوصول بعد الإلغاء، مهما كان تسلسلها في الشجرة \\
\bottomrule
\end{tabular}
\end{center}

\subsection{جدولة \en{MLFQ} --- تصنيف A}

خوارزمية الطوابير المتعددة بتغذية رجعية (Multi-Level Feedback Queue) تدير 4 طوابير بأوزان زمنية متفاوتة، وتحقق توازناً مثالياً بين الاستجابية وإنصاف العمليات الطويلة.

\vspace{0.3cm}

\begin{center}
\begin{tikzpicture}[
  queue/.style={draw=accent, fill=accent!12, rounded corners=3pt,
    minimum width=7cm, minimum height=0.8cm, align=left},
  proc/.style={draw=accent2, fill=accent2!25, rounded corners=2pt,
    minimum width=1.0cm, minimum height=0.6cm, font=\tiny},
  timebox/.style={draw=accentg, fill=accentg!15, rounded corners=2pt,
    minimum width=1.2cm, minimum height=0.6cm, font=\small\bfseries}
]

\node[queue, label=left:{\small\bfseries \en{Q0} (أعلى)}] (q0) at (0,3) {};
\node[queue, label=left:{\small\bfseries \en{Q1} (عالية)}] (q1) at (0,2) {};
\node[queue, label=left:{\small\bfseries \en{Q2} (متوسطة)}] (q2) at (0,1) {};
\node[queue, label=left:{\small\bfseries \en{Q3} (منخفضة)}] (q3) at (0,0) {};

% وقت الكم
\node[timebox] (t0) at (5.5,3) {8 ms};
\node[timebox] (t1) at (5.5,2) {16 ms};
\node[timebox] (t2) at (5.5,1) {32 ms};
\node[timebox] (t3) at (5.5,0) {64 ms};

% عمليات في الطوابير
\foreach \x/\c in {-3.2/accent!50, -2.0/accent!40, -0.8/accent!30}{
  \node[proc, fill=\c] at (\x,3) {P};
}
\foreach \x/\c in {-3.2/accent2!40, -2.0/accent2!30}{
  \node[proc, fill=\c] at (\x,2) {P};
}
\node[proc] at (-3.2,1) {P};

% أسهم التنزيل
\draw[->, dashed, accent3, thick]
  (3.5,2.8) to[bend right=20]
  node[right, font=\scriptsize, text=accent3]
    {تجاوز القانتوم $\rightarrow$ انتقال للأسفل} (3.5,2.1);
\draw[->, dashed, accent3, thick]
  (3.5,1.8) to[bend right=20] (3.5,1.1);

% رفع الأولوية كل 1000ms
\draw[->, accent2, very thick, bend left=50]
  (-3.8,-0.2) to
  node[left, font=\scriptsize, text=accent2]
    {رفع كل 1000ms} (-3.8,2.8);

\node[font=\footnotesize, text=darkgray] at (0,-0.8)
  {جدولة MLFQ --- 4 طوابير، أولوية ديناميكية، منع التجويع};
\end{tikzpicture}
\captionof{figure}{نظام جدولة \en{MLFQ} مع 4 طوابير وآلية رفع الأولوية}
\end{center}

\vspace{0.3cm}

نموذج الأولوية الرسمي:

\begin{equation}
\mathrm{Priority}(p) = \begin{cases}
\mathrm{Q0} & \text{if } \mathrm{yield\_count}(p) < 5 \;\lor\; \mathrm{boost\_flag} = \mathrm{true}\\
\mathrm{Q1} & \text{if } \mathrm{voluntary\_ctx}(p) > \mathrm{involuntary\_ctx}(p)\\
\mathrm{Q2} & \text{if } \mathrm{cpu\_time}(p) < 100\;\mathrm{ms}\\
\mathrm{Q3} & \text{otherwise}
\end{cases}
\end{equation}

رفع الأولوية الدوري كل $T_{\mathrm{boost}} = 1000\;\mathrm{ms}$ لمنع التجويع:
\begin{equation}
\forall\, p \in \mathrm{Processes}:\; \mathrm{Priority}(p) \gets \mathrm{Q0}
\end{equation}

\vspace{0.3cm}

\begin{center}
\begin{tabular}{C{1.5cm}C{1.5cm}C{1.5cm}L{6cm}}
\toprule
\rowcolor{titlebg!85}
\bfseries\color{white}\textarabic{الطابور} &
\bfseries\color{white}\textarabic{الأولوية} &
\bfseries\color{white}\textarabic{الكم الزمني} &
\bfseries\color{white}\textarabic{الاستخدام المثالي} \\
\midrule
\rowcolor{rowalt} \en{Q0} & أعلى & 8 \en{ms} & عمليات تفاعلية (\en{Shell}, \en{Editor}) وجميع العمليات الجديدة \\
\en{Q1} & عالية & 16 \en{ms} & عمليات \en{CPU-bound} طفيفة الاستخدام \\
\rowcolor{rowalt} \en{Q2} & متوسطة & 32 \en{ms} & عمليات حسابية طويلة \\
\en{Q3} & منخفضة & 64 \en{ms} & مهام خلفية (تجميع، نسخ احتياطي) \\
\bottomrule
\end{tabular}
\end{center}

\subsection{تحصين المجدول ومقاومة المقاطعات المتزامنة (\en{Scheduler Hardening})}

لضمان سلامة النواة تحت الأعباء العالية وتعدد المعالجات (\en{SMP})، تم إدخال تحسينات حاسمة على نظام الجدولة وهيكل النواة الأساسي:

\begin{itemize}[leftmargin=2em]
\item \textbf{منع التعليق المتبادل (\en{Deadlock Avoidance}):}
تم تغليف عمليات المجدول الحيوية (مثل \codeinline{spawn}، \codeinline{yield\_now}، \codeinline{wake\_sleeping}، \codeinline{exec\_process}) بكتل حظر مقاطعات كاملة عبر \codeinline{x86\_64::instructions::interrupts::without\_interrupts}. يحمي هذا الإجراء أقفال المجدول المشتركة من التعليق المتبادل عندما تنطلق مقاطعة المؤقت الدوري (\en{Timer Interrupt}) وتطلب قفل الجدولة بينما يحتفظ به معالج آخر في نفس اللحظة.

\item \textbf{أجهزة تتبع الإقلاع المبكر (\en{UART Raw Logging}):}
تم تأسيس آلية طباعة تسلسلية مباشرة للعتاد \codeinline{raw\_serial\_log} للعمل بدون ذاكرة مخصصة (\en{no\_std early boot})، مما سمح بمراقبة دقيقة وخطوة بخطوة لمراحل إعداد الـ \en{GDT}، \en{IDT}، \en{PIC}، والذاكرة السطحية وهيكل APIC، لتسهيل تشخيص أعطال التعليق المبكر قبل تهيئة نظام الطباعة القياسي.

\item \textbf{إعادة تنظيم تهيئة الأمان (\en{Security Feature Ordering}):}
أعيد ترتيب خطوات تمكين ميزات الحماية الفيزيائية للمنافذ والذاكرة (\en{SMEP, SMAP, UMIP}) عبر سجل التحكم \codeinline{CR4}. حيث أرجئ تفعيلها إلى نهاية مرحلة الإقلاع التام لضمان تهيئة ذاكرة الكومة (\en{Heap}) وشاشة العرض ونظام التقاط الأخطاء التسلسلي أولاً، متفادين بذلك حدوث انهيار عام (\en{Kernel Panic}) فوري عند قراءة صفحات الذاكرة العشوائية التابعة للمستخدم.

\item \textbf{إصلاح أزاحات منافذ الشبكة (\en{VirtIO-Net Correction}):}
تم تصحيح إزاحات واجهة منافذ PCI للبطاقة الافتراضية (\en{VirtIO Network Virtual Device}) بتعديل إزاحة مسجل مقاطعة العتاد \codeinline{PCI\_ISR} إلى \codeinline{0x13} ومسجل إعدادات العتاد \codeinline{PCI\_DEVICE\_CFG} إلى \codeinline{0x14}، مما سمح للنواة باستقبال مقاطعات بطاقة الشبكة ومعالجة البيانات الواردة بدقة.

\item \textbf{علاج ضغط مكدس النواة (\en{Stack Pressure Relief}):}
بسبب طبيعة الموارد المحدودة لمكدس خيط النواة (\en{Kernel Stack}) البالغ 64 كيلوبايت، تم نقل كائنات واجهة الشبكة ومجموعة المقابس التابعة لمكتبة \codeinline{smoltcp} من التخصيص الداخلي في المكدس إلى التخصيص الديناميكي على الكومة (\en{Heap-allocation} عبر \codeinline{Box})، مما أزال خطر انهيار مكدس النواة (\en{Stack Overflow}) تماماً عند معالجة الحزم المتتالية.
\end{itemize}

\subsection{إدارة الذاكرة}

\begin{center}
\begin{tikzpicture}[scale=0.85]
% مخطط توزيع الذاكرة الافتراضية
\def\W{6}
\def\startY{0}

% Boot Info
\fill[orange!25, draw=orange!60] (0,\startY) rectangle (\W, \startY+0.5);
\node[font=\tiny] at (\W/2, \startY+0.25) {Boot Info \texttt{0x0--0xFFFF}};

% Kernel Code
\fill[accent!25, draw=accent!60] (0,\startY+0.5) rectangle (\W, \startY+1.4);
\node[font=\tiny] at (\W/2, \startY+0.95) {Kernel Code \texttt{0x100000--0x200000}};

% Kernel Stack
\fill[accent!40, draw=accent!70] (0,\startY+1.4) rectangle (\W, \startY+2.0);
\node[font=\tiny] at (\W/2, \startY+1.7) {Kernel Stack};

% Page Tables
\fill[accent2!25, draw=accent2!60] (0,\startY+2.0) rectangle (\W, \startY+2.8);
\node[font=\tiny] at (\W/2, \startY+2.4) {Page Tables (L4 Shared)};

% Heap
\fill[teal!25, draw=teal!60] (0,\startY+2.8) rectangle (\W, \startY+4.0);
\node[font=\small] at (\W/2, \startY+3.4) {Kernel Heap};

% فراغ
\fill[gray!10, draw=gray!30, dashed] (0,\startY+4.0) rectangle (\W, \startY+5.0);
\node[font=\tiny, text=gray] at (\W/2, \startY+4.5) {...};

% User Space
\fill[accent2!20, draw=accent2!50] (0,\startY+5.0) rectangle (\W, \startY+6.5);
\node[font=\small] at (\W/2, \startY+5.75) {User Space \texttt{0x300000+}};

% أسهم العناوين
\node[font=\tiny, text=darkgray, anchor=east] at (-0.1, \startY+0.25) {\texttt{0x0}};
\node[font=\tiny, text=darkgray, anchor=east] at (-0.1, \startY+0.5) {\texttt{0x10000}};
\node[font=\tiny, text=darkgray, anchor=east] at (-0.1, \startY+1.4) {\texttt{0x200000}};
\node[font=\tiny, text=darkgray, anchor=east] at (-0.1, \startY+5.0) {\texttt{0x300000+}};

% تسمية
\node[font=\small\bfseries, text=darkgray, anchor=west] at (\W+0.3, \startY+3.25)
  {فضاء العناوين الافتراضية};
\draw[dashed, gray] (\W, \startY) -- (\W+0.3, \startY);
\draw[dashed, gray] (\W, \startY+6.5) -- (\W+0.3, \startY+6.5);
\draw[<->, gray] (\W+0.2, \startY) -- (\W+0.2, \startY+6.5);
\end{tikzpicture}
\captionof{figure}{تخطيط الذاكرة الافتراضية لـ \textsc{ZiqaKernel}}
\end{center}

\vspace{0.3cm}

\subsubsection{ذاكرة التخزين المؤقت LRU}

نسبة الإصابة (Hit Rate) تصل إلى \en{85.4\%} مع وصول O(1) عبر قاموس تجزئة وقائمة LRU:

\begin{equation}
\mathrm{HitRate} = \frac{H}{H + M} \approx 0.854,\qquad
H = \text{عدد الإصابات},\quad M = \text{عدد التخطيات}
\end{equation}

\begin{LTR}
\begin{lstlisting}[style=rustcode, caption={تنفيذ LRU Page Cache بوصول O(1)}]
struct PageCache {
    entries:  HashMap<u64, PageEntry>,
    lru_list: VecDeque<u64>,
    hits:     AtomicU64,
    misses:   AtomicU64,
}

fn access(&mut self, page_id: u64) -> Option<&PageEntry> {
    if self.entries.contains_key(&page_id) {
        // (*@نقل الصفحة للمقدمة (مستخدمة مؤخراً)@*)
        self.lru_list.retain(|x| *x != page_id);
        self.lru_list.push_front(page_id);
        self.hits.fetch_add(1, Ordering::Relaxed);
        self.entries.get(&page_id)
    } else {
        self.misses.fetch_add(1, Ordering::Relaxed);
        None
    }
}
\end{lstlisting}
\end{LTR}

\newpage

% ══════════════════════════════════════════════════════════════════
\section{نظام eBPF --- التحقق والتنفيذ}\label{sec:ebpf}
% ══════════════════════════════════════════════════════════════════

نظام \textsc{eBPF} (Extended Berkeley Packet Filter) في \textsc{ZiqaKernel} يُصنَّف \en{Obsidian-Tier}. يتكون من 2,583 سطر موزعة على 5 وحدات ويوفر بيئة آمنة لتنفيذ برامج المستخدم داخل النواة.

\subsection{مكونات النظام الخمس}

\begin{center}
\begin{tikzpicture}[
  mod/.style={draw=accent4, fill=accent4!12, rounded corners=5pt,
    minimum width=2.5cm, minimum height=1.4cm,
    text width=2.2cm, align=center, font=\small},
  arrow/.style={->, thick, accent4, shorten >=2pt, shorten <=2pt}
]
\node[mod] (vm) at (0,0) {\texttt{vm.rs}\\[3pt]637 سطر\\\tiny المفسّر};
\node[mod, right=1.2cm of vm] (ver)
  {\texttt{verifier.rs}\\[3pt]1076 سطر\\\tiny المدقق};
\node[mod, right=1.2cm of ver] (map)
  {\texttt{map.rs}\\[3pt]187 سطر\\\tiny الخرائط};
\node[mod, below=1.2cm of vm] (att)
  {\texttt{attach.rs}\\[3pt]128 سطر\\\tiny الربط};
\node[mod, right=1.2cm of att] (mod_)
  {\texttt{mod.rs}\\[3pt]555 سطر\\\tiny الأنواع};

\draw[arrow] (ver) -- (vm) node[midway,above,font=\tiny] {إجازة};
\draw[arrow] (att) -- (vm) node[midway,left,font=\tiny] {إطلاق};
\draw[arrow] (vm) -- (map) node[midway,above,font=\tiny,sloped] {وصول};
\draw[arrow] (mod_) -- (ver) node[midway,above,font=\tiny] {أنواع};
\draw[arrow] (mod_) -- (vm) node[midway,right,font=\tiny] {تعليمات};
\end{tikzpicture}
\captionof{figure}{وحدات نظام eBPF وتفاعلاتها}
\end{center}

\subsection{معمارية الآلة الافتراضية}

الآلة الافتراضية تحاكي معالجاً بـ 11 سجلاً (R0--R10) ومكدس ثابت 512 بايت. كل تعليمة 64-بت:

\begin{equation}
\underbrace{\texttt{opcode}}_{8\;\text{bits}} \;\Big|\;
\underbrace{\texttt{dst\_src}}_{4+4\;\text{bits}} \;\Big|\;
\underbrace{\texttt{offset}}_{16\;\text{bits}} \;\Big|\;
\underbrace{\texttt{immediate}}_{32\;\text{bits}}
\end{equation}

\begin{center}
\begin{tabular}{C{2.5cm}L{5cm}L{4.5cm}}
\toprule
\rowcolor{accent4!20}
\bfseries\color{accent4}\textarabic{الفئة} &
\bfseries\color{accent4}\textarabic{التعليمات} &
\bfseries\color{accent4}\textarabic{الوصف} \\
\midrule
\rowcolor{rowalt} ALU64 &
  ADD, SUB, MUL, DIV, OR, AND,\ldots &
  عمليات 64-بت على السجلات \\
ALU32 &
  MOV32, ADD32, SUB32,\ldots &
  عمليات 32-بت مع تمديد صفري \\
\rowcolor{rowalt} Memory &
  LDX, STX (8/16/32/64-bit) &
  تحميل وتخزين الذاكرة \\
Jump &
  JA, JEQ, JGT, JLT, JNE,\ldots &
  قفزات شرطية وغير شرطية \\
\rowcolor{rowalt} CALL/EXIT &
  CALL (helper), EXIT &
  استدعاء المساعدين والخروج \\
\bottomrule
\end{tabular}
\end{center}

\subsection{مدقق التحليل التجريدي}

المدقق هو أعقد مكون في النظام. يستخدم التحليل التجريدي (\en{Abstract Interpretation}) لضمان سلامة كل برنامج eBPF قبل تنفيذه.

\begin{definitionbox}[الحالة المجردة]
في كل نقطة تنفيذ، تُمثَّل حالة كل سجل بثلاثة عناصر:
\begin{equation}
\sigma : \{R_0, \ldots, R_{10}\} \to \underbrace{\mathrm{RegType}}_{\text{6 أنواع}}
  \times \underbrace{[\min,\max]}_{\text{حدود القيمة}}
  \times \underbrace{\mathrm{ZeroUpper}}_{\text{bit 32 حالة}}
\end{equation}
حيث أنواع السجلات:
\codeinline{Uninitialized | Scalar | PtrToStack | PtrToMapValue | PtrToCtx | PtrToMapKey}
\end{definitionbox}

\begin{center}
\begin{tikzpicture}[
  stepnode/.style={draw=accent4, fill=accent4!10, rounded corners=5pt,
    minimum width=3.5cm, minimum height=1.0cm, align=center, font=\small},
  arrow/.style={->, thick, accent4}
]
\node[stepnode] (s1) at (0,0) {بناء CFG\\\tiny Control Flow Graph};
\node[stepnode] (s2) at (4.5,0) {التفسير التجريدي\\\tiny Abstract Interpretation};
\node[stepnode] (s3) at (9,0) {التحقق من الذاكرة\\\tiny Memory Access Check};
\node[stepnode, below=1cm of s2] (s4) {نقطة الثبات\\\tiny Fixpoint Convergence};
\node[draw=accent2, fill=accent2!10, rounded corners=5pt,
      right=1cm of s3, minimum height=1cm] (ok) {\color{accent2}\bfseries OK};
\node[draw=accent3, fill=accent3!10, rounded corners=5pt,
      below right=0.5cm and 1cm of s3, minimum height=1cm] (err)
  {\color{accent3}\bfseries ERROR};

\draw[arrow] (s1) -- (s2);
\draw[arrow] (s2) -- (s3);
\draw[arrow] (s2) -- (s4);
\draw[arrow] (s4) to[bend right=30] (s2);
\draw[arrow] (s3) -- (ok);
\draw[arrow] (s3) -- (err);
\end{tikzpicture}
\captionof{figure}{خط أنابيب مدقق eBPF بالتحليل التجريدي}
\end{center}

\begin{theorem}[إنهاء المدقق]
يضمن مدقق \textsc{ZiqaKernel} الإنهاء في زمن محدود عبر ثلاث آليات:
\begin{enumerate}[leftmargin=2em]
\item شبكة محدودة من أنواع السجلات ($|\mathrm{RegType}| = 6$ $\rightarrow$ ضمان التقارب)
\item دمج تحفظي عند نقاط الالتقاء (merge $\to$ Scalar عند الاختلاف)
\item حد أقصى للتعليمات $N_{\max} = 100{,}000$ لمنع الحلقات اللانهائية
\end{enumerate}
\end{theorem}

\vspace{0.4cm}

دالة الدمج عند نقاط الالتقاء:
\begin{equation}
\mathrm{merge}(\sigma_1, \sigma_2) = \begin{cases}
t_1 & \text{if } t_1 = t_2 \\
\texttt{Scalar} & \text{otherwise}
\end{cases}
\qquad
\mathrm{merge}([a,b],\,[c,d]) = [\min(a,c),\, \max(b,d)]
\end{equation}

\subsection{ميزات \textsc{eBPF} المتقدمة (أطروحة \en{Obsidian-Tier})}

شهد نظام \textsc{eBPF} تحديثات جوهرية نقلته إلى مرحلة نضوج كاملة، بهدف تمكين التطبيقات الأكثر تعقيداً داخل النواة:

\begin{enumerate}[leftmargin=2em]
\item \textbf{الحلقات المقيدة (\en{Bounded Loops}):} 
تمت ترقية المدقق (\codeinline{BpfVerifier}) للتخلي عن منع القفزات الخلفية المطلق، واستبداله بتحليل ديناميكي يسمح بالحلقات الخلفية المقيدة بحدود واضحة ومعلومة أثناء التحقق التجريدي، مع تعيين حد أقصى للتعليمات لتفادي الحلقات اللانهائية والـ \en{Denial of Service}.

\item \textbf{خرائط التجزئة والبرمجيات (\en{Hash Maps \& ProgArray}):}
تمت إضافة نوع خريطة التجزئة \codeinline{BpfMapType::Hash} باستخدام التجزئة المفتوحة وعلاج التصادمات بالاستقصاء الخطي (\en{Linear Probing})، بالإضافة إلى خريطة برمجيات \codeinline{ProgArray} التي تُمكِّن الاستدعاءات الذيلية (\en{Tail Calls}) عبر مساعد \codeinline{bpf\_tail\_call} لتغيير سياق التنفيذ لبرنامج آخر ديناميكياً بدون استهلاك مكدس النواة.

\item \textbf{المساعدون الجدد لسلامة الوصول والأنظمة (\en{Kernel Helpers}):}
\begin{itemize}
\item \codeinline{bpf\_probe\_read}: لقراءة الذاكرة العشوائية بأمان تام مع معالجة أخطاء الصفحة تلقائياً.
\item \codeinline{bpf\_get\_current\_comm}: لقراءة اسم العملية الحالية في فضاء المستخدم وحفظه في مصفوفة محددة.
\item \codeinline{GET\_SMP\_PROCESSOR\_ID}: مساعد متوافق مع النوى المتعددة للاستعلام عن معرّف المعالج الفيزيائي الحالي للجدولة المتوازنة.
\end{itemize}

\item \textbf{مكدس الآلة الافتراضية والتعليمات الموسعة (\en{VM Stack \& Memory Insts}):}
دعم كامل للذاكرة المستقلة للآلة الافتراضية عبر تعليمات التحميل والتخزين غير المباشر (\codeinline{LDX, STX, ST}) وربطها بمكدس تنفيذي خاص بكل آلة لتوفير مساحة لتخزين المتغيرات المحلية.
\end{enumerate}

\newpage

% ══════════════════════════════════════════════════════════════════
\section{تحليل الأداء}\label{sec:perf}
% ══════════════════════════════════════════════════════════════════

\subsection{مقاييس الأداء الرئيسية}

\begin{center}
\begin{tabular}{L{4cm}C{3cm}C{3cm}C{3cm}}
\toprule
\rowcolor{titlebg!85}
\bfseries\color{white}\textarabic{المقياس} &
\bfseries\color{white}\textarabic{ZiqaKernel} &
\bfseries\color{white}\textarabic{Linux 6.x} &
\bfseries\color{white}\textarabic{seL4} \\
\midrule
\rowcolor{rowalt} Boot Time & $< 1\;\mathrm{s}$ & $10$--$30\;\mathrm{s}$ & $< 0.1\;\mathrm{s}$ \\
Context Switch & $< 5\;\mu\mathrm{s}$ & $< 10\;\mu\mathrm{s}$ & $< 1\;\mu\mathrm{s}$ \\
\rowcolor{rowalt} Heap Efficiency & \en{98.2\%} & $\sim$\en{95\%} & $\sim$\en{99\%} \\
Cache Hit Rate & \en{85.4\%} & $\sim$\en{80\%} & N/A \\
\rowcolor{rowalt} استهلاك الذاكرة & $\sim$2 \en{MB} & $\sim$5 \en{MB} & $< 1\;\mathrm{MB}$ \\
\en{Network} Throughput & 5.2 \en{GB/s} & $\sim$10 \en{GB/s} & N/A \\
\rowcolor{rowalt} Storage Throughput & 8.7 \en{GB/s} & $\sim$5 \en{GB/s} & N/A \\
سطور الكود & 28,633 & $\sim$28M & $\sim$9K \\
\bottomrule
\end{tabular}
\end{center}

\subsection{قانون أمدال --- التحليل الرياضي}

مع نسبة توازي $P = 0.75$ (\en{75\%} من الكود قابل للتوزيع)، نتوقع وفق قانون أمدال:

\begin{equation}
\mathrm{Speedup}(N) = \frac{1}{(1-P) + \dfrac{P}{N}},\qquad P = 0.75,\; N = \text{عدد النوى}
\end{equation}

\begin{center}
\begin{tikzpicture}
\begin{axis}[
  x dir=reverse,
  width=0.85\textwidth, height=7cm,
  xlabel={عدد النوى $N$},
  ylabel={معامل التسريع},
  xmin=1, xmax=32, ymin=1, ymax=35,
  xtick={1,2,4,8,16,32},
  ytick={1,4,7,10,15,20,25,30},
  legend pos=north west,
  ymajorgrids=true, xmajorgrids=true,
  grid style={dashed, gray!25},
  axis line style={gray!60},
]
% أمدال الفعلي
\addplot[color=accent2, thick, mark=*, mark size=3pt,
  mark options={fill=accent2}] coordinates {
  (1,1) (2,1.6) (4,2.29) (8,3.2) (16,4.0) (32,4.57)
};
\addlegendentry{Amdahl ($P=0.75$)}
% مثالي خطي
\addplot[color=accent, dashed, thick] coordinates {
  (1,1) (2,2) (4,4) (8,8) (16,16) (32,32)
};
\addlegendentry{مثالي (خطي)}
% أمدال بتوازي 90%
\addplot[color=orange, dotted, thick, mark=square*, mark size=2pt,
  mark options={fill=orange}] coordinates {
  (1,1) (2,1.82) (4,3.07) (8,4.71) (16,6.4) (32,7.8)
};
\addlegendentry{Amdahl ($P=0.90$)}

\node[font=\tiny, text=darkgray] at (axis cs:16,3.5)
  {المتوقع: 3.6× على 8 نوى};
\end{axis}
\end{tikzpicture}
\captionof{figure}{قانون أمدال: تسريع متعدد النوى لمختلف نسب التوازي}
\end{center}

\subsection{تقييم الأداء متعدد الأبعاد}

\begin{center}
\begin{tikzpicture}
\begin{axis}[
  x dir=reverse,
  ybar,
  width=0.9\textwidth, height=6.5cm,
  symbolic x coords={الإقلاع,تبديل السياق,كفاءة الذاكرة,
    أمان القدرات,نظام eBPF,التوافق},
  xtick=data,
  xticklabel style={font=\small, rotate=25, anchor=east},
  ylabel={التقييم (\%)},
  ymin=0, ymax=105,
  nodes near coords,
  nodes near coords align={vertical},
  every node near coord/.append style={font=\small, color=white,
    rotate=90, anchor=west, yshift=4pt},
  bar width=16pt,
  ymajorgrids=true,
  grid style={dashed, gray!25},
  legend style={at={(0.97,0.97)}, anchor=north east, font=\small},
]
\addplot[fill=accent!80, draw=accent, fill opacity=0.85] coordinates {
  (الإقلاع,95)
  (تبديل السياق,90)
  (كفاءة الذاكرة,98)
  (أمان القدرات,99)
  (نظام eBPF,92)
  (التوافق,77)
};
\addlegendentry{ZiqaKernel}

\addplot[fill=accent2!70, draw=accent2, fill opacity=0.85] coordinates {
  (الإقلاع,50)
  (تبديل السياق,70)
  (كفاءة الذاكرة,80)
  (أمان القدرات,60)
  (نظام eBPF,95)
  (التوافق,100)
};
\addlegendentry{Linux 6.x}

\addplot[fill=orange!70, draw=orange, fill opacity=0.85] coordinates {
  (الإقلاع,99)
  (تبديل السياق,99)
  (كفاءة الذاكرة,99)
  (أمان القدرات,100)
  (نظام eBPF,30)
  (التوافق,40)
};
\addlegendentry{seL4}
\end{axis}
\end{tikzpicture}
\captionof{figure}{مقارنة الأداء متعدد الأبعاد بين الأنوية الثلاث}
\end{center}

\subsection{أداء الكم الزمني وكفاءة الأقفال}

مع fine-grained locking وعند $P_{\mathrm{contention}} = 0.1$ و $T_{\mathrm{lock}} = 0.1\;\mu\mathrm{s}$:

\begin{equation}
T_{\mathrm{effective}} = T_{\mathrm{compute}} + T_{\mathrm{lock}}
  \times \frac{N(N-1)}{2} \times P_{\mathrm{contention}}
\end{equation}

يبقى تأثير القفل أقل من \en{5\%} حتى 8 أنوية، مما يُثبت كفاءة نموذج fine-grained locking.

\vspace{0.3cm}

% رسم 3D للأداء
\begin{center}
\begin{tikzpicture}
\begin{axis}[
  x dir=reverse,
  view={40}{30},
  width=0.9\textwidth, height=8cm,
  xlabel={النوى},
  ylabel={العمليات/ثانية (K)},
  zlabel={الكفاءة (\%)},
  xtick={1,2,4,8},
  grid=both,
  grid style={line width=0.3pt, gray!30},
  colormap/viridis,
  colorbar,
  colorbar style={ylabel={الكفاءة}},
  zmin=50, zmax=100,
  title={العلاقة ثلاثية الأبعاد: النوى × الإنتاجية × الكفاءة},
]
\addplot3[surf, shader=faceted interp, samples=8, domain=1:8,
  y domain=100:800, opacity=0.9, faceted color=black!60] {
  100*(0.75*(x-1)/x + 0.25)*(y/800)*
  (1 - 0.02*(x-1))
};
\end{axis}
\end{tikzpicture}
\captionof{figure}{رسم ثلاثي الأبعاد للعلاقة بين عدد النوى والإنتاجية والكفاءة}
\end{center}

\newpage

% ══════════════════════════════════════════════════════════════════
\section{مقارنة مع أنوية أخرى}\label{sec:compare}
% ══════════════════════════════════════════════════════════════════

\subsection{فلسفة التصميم المقارنة}

\begin{center}
\begin{tabular}{L{3cm}C{3.3cm}C{3.3cm}C{3.3cm}}
\toprule
\rowcolor{titlebg!85}
\bfseries\color{white}\textarabic{البُعد} &
\bfseries\color{white}\textarabic{ZiqaKernel} &
\bfseries\color{white}\textarabic{Linux 6.x} &
\bfseries\color{white}\textarabic{seL4} \\
\color{white}\bfseries seL4 \\
\midrule
\rowcolor{rowalt} اللغة & Rust + Zig & C & C + ASM \\
نوع النواة & \en{Microkernel} & Monolithic & \en{Microkernel} \\
\rowcolor{rowalt} نموذج الأمان & Capability & DAC/MAC/AppArmor & متحقق رسمياً \\
التحقق الرسمي & جزئي & غير موجود & كامل (Isabelle) \\
\rowcolor{rowalt} عدد الملفات & 101 & $\sim$30,000 & $\sim$100 \\
سطور الكود & 28,633 & $\sim$28M & $\sim$9K \\
\rowcolor{rowalt} استدعاءات النظام & 129 & 400+ & 50+ \\
الجدولة & \en{MLFQ} (4 طوابير) & CFS & Round-Robin \\
\rowcolor{rowalt} eBPF & \en{Obsidian-Tier} & Mature (BTF) & غير متاح \\
برمجة ديناميكية & eBPF (كاملة) & eBPF + kprobes & محدودة \\
\bottomrule
\end{tabular}
\end{center}

\subsection{مقارنة نظام eBPF}

\begin{center}
\begin{tabular}{L{3.5cm}L{4.5cm}L{4.5cm}}
\toprule
\rowcolor{accent4!20}
\bfseries\color{accent4}\textarabic{الخاصية} &
\bfseries\color{accent4}\textarabic{ZiqaKernel eBPF} &
\bfseries\color{accent4}\textarabic{Linux eBPF} \\
\color{accent4}\bfseries \en{Linux} eBPF \\
\midrule
\rowcolor{rowalt} حجم المكدس & 512B ثابت & 512B قابل للتعديل \\
أنواع الخرائط & Array + Hash + 3 أخرى & 12+ نوعاً (ringbuf, lru, \ldots) \\
\rowcolor{rowalt} المساعدون & 5 أساسيون & 100+ \\
المدقق & CFI + kCFI + تحليل تجريدي & CFI + BTF Bounds \\
\rowcolor{rowalt} تتبع الأنواع & كامل (6 أنواع) & كامل (BTF) \\
اللغة & Rust (آمن) & C \\
\rowcolor{rowalt} التحقق الرياضي & تحليل تجريدي & تحليل تدفق \\
دعم SMP & نعم & نعم \\
\bottomrule
\end{tabular}
\end{center}

\subsection{الاختلافات الجوهرية}

\begin{resultbox}[\en{ZiqaKernel} مقابل \en{Linux}]
\textbf{ZiqaKernel} يُفضِّل السلامة عبر Rust والقدرات على الأداء الخام. مساحة السطح أصغر بعامل 1000× (28K مقابل 28M سطر)، مما يجعل التحقق الرسمي ممكناً. يفتقر إلى نضج سائقي الأجهزة ولكنه يتفوق في نموذج الأمان وسرعة الإقلاع.
\end{resultbox}

\begin{resultbox}[\en{ZiqaKernel} مقابل seL4]
\textbf{ZiqaKernel} يقدم مرونة أكبر عبر eBPF الديناميكي وجدولة \en{MLFQ} متكيفة على حساب غياب التحقق الرسمي الكامل. كلاهما يستخدم نموذج القدرات، لكن \en{ZiqaKernel} مكتوب بـ Rust (آمن بالتصميم) بدلاً من C.
\end{resultbox}

\newpage

% ══════════════════════════════════════════════════════════════════
\section{تحليل خريطة المعرفة}\label{sec:graph}
% ══════════════════════════════════════════════════════════════════

\subsection{إحصائيات الرسم البياني المعرفي}

\begin{center}
\begin{tabular}{L{5cm}R{4cm}}
\toprule
\rowcolor{titlebg!85}
\bfseries\color{white}\textarabic{المقياس} & \bfseries\color{white}\textarabic{القيمة} \\
\midrule
\rowcolor{rowalt} إجمالي العقد & 2,015 \\
إجمالي العلاقات & 3,843 \\
\rowcolor{rowalt} عدد المجتمعات & 131 \\
عقد الكود (Code Symbols) & 1,660 \\
\rowcolor{rowalt} عقد الوثائق & 347 \\
ملفات المصدر & 101 \\
\rowcolor{rowalt} متوسط الدرجة & 3.82 \\
معامل التجميع & 0.34 \\
\rowcolor{rowalt} الطول الوسطى للمسار & 4.7 حافة \\
توفير الرمز المميز (\en{graphify}) & 75.4× \\
\bottomrule
\end{tabular}
\end{center}

\subsection{توزيع المجتمعات الخمسة عشر الأولى}

\begin{center}
\begin{tikzpicture}
\begin{axis}[
  y dir=reverse,
  xbar,
  width=\textwidth, height=9cm,
  symbolic y coords={
    Build,Cap,IPC,Net,Syscall,VGA,\en{ABI},Mem,Editor,Win,Sched.,PCI,Shell,\en{ABI}-L,ZiqaFS},
  ytick=data,
  yticklabel style={font=\small},
  xlabel={عدد العقد},
  xmin=0, xmax=105,
  nodes near coords,
  nodes near coords align={horizontal},
  every node near coord/.append style={font=\footnotesize},
  bar width=10pt,
  enlarge y limits=0.05,
  ymajorgrids=true,
  grid style={dashed, gray!25},
]
\addplot[fill=accent!70, draw=accent, fill opacity=0.85] coordinates {
  (34,Build) (35,Cap) (36,IPC) (37,Net) (38,Syscall) (39,VGA) (39,\en{ABI})
  (53,Mem) (59,Editor) (72,Win) (80,Sched.) (84,PCI) (86,Shell) (87,\en{ABI}-L) (90,ZiqaFS)
};
\end{axis}
\end{tikzpicture}
\captionof{figure}{توزيع العقد عبر المجتمعات المعمارية الخمسة عشر الأولى}
\end{center}

\subsection{اكتشاف المجتمعات بخوارزمية لوفان}

خوارزمية لوفان (\en{Louvain}) تُعظِّم معامل النمطية (Modularity):

\begin{equation}
Q = \frac{1}{2m}\sum_{i,j}\left[A_{ij} - \frac{k_i k_j}{2m}\right]\delta(c_i, c_j)
\end{equation}

حيث: $m$ = عدد الحواف الكلي، $k_i$ = درجة العقدة $i$، $c_i$ = مجتمع العقدة $i$، $\delta$ = دالة كرونيكر دلتا.

\newpage

% ══════════════════════════════════════════════════════════════════
\section{الأبحاث المتقدمة والتحديات}
% ══════════════════════════════════════════════════════════════════

\subsection{تحديات التزامن في المتعدد النوى}

أحد أصعب المشكلات في تطوير الأنوية متعددة النوى هو التزامن. عند 8 أنوية، قد تتنافس عمليات متعددة على نفس المورد في نفس الوقت.

حالة السباق (Race Condition) في استدعاءات النظام:

\begin{equation}
\mathrm{Latency}_{\mathrm{lock}} = \frac{N_{\mathrm{ops}} \times T_{\mathrm{atomic}}}{N_{\mathrm{cores}}}
  + T_{\mathrm{critical}} \times P_{\mathrm{conflict}}
\end{equation}

الحل المُنفَّذ: Fine-grained locking مع قفل منفصل لكل:
\begin{itemize}[leftmargin=2em]
\item قناة IPC (per-channel locks) ---يُتيح IPC موازياً كاملاً
\item جدول VFS لكل عملية (per-process VFS table)
\item RwLock على جدولة العمليات (قراءة متوازية، كتابة حصرية)
\end{itemize}

\subsection{أمان الذاكرة: SMEP / SMAP / UMIP}

حماية ثلاثية الطبقات ضد هجمات فضاء المستخدم:

\begin{center}
\begin{tikzpicture}[
  layer/.style={draw, rounded corners=4pt, minimum width=9cm,
    minimum height=0.8cm, font=\small, align=center},
  anno/.style={font=\tiny, text=darkgray}
]
\node[layer, fill=accent3!15, draw=accent3] (smep) at (0,2)
  {SMEP (CR4.20) --- منع النواة من تنفيذ كود المستخدم};
\node[layer, fill=orange!15, draw=orange] (smap) at (0,1)
  {SMAP (CR4.21) --- منع النواة من الوصول لذاكرة المستخدم بدون STAC/CLAC};
\node[layer, fill=accent4!15, draw=accent4] (umip) at (0,0)
  {UMIP (CR4.11) --- منع المستخدم من تنفيذ SGDT/SIDT/SMSW};

\draw[<->, accent3, dashed] (smep.east) -- ++(1,0)
  node[right, anno] {Ring 3 $\rightarrow$ Ring 0 انتهاك};
\draw[<->, orange, dashed] (smap.east) -- ++(1,0)
  node[right, anno] {STAC/CLAC تحكم};
\draw[<->, accent4, dashed] (umip.east) -- ++(1,0)
  node[right, anno] {منع تسريب GDTR/IDTR};
\end{tikzpicture}
\captionof{figure}{الطبقات الثلاث للحماية الأمنية للذاكرة}
\end{center}

\begin{LTR}
\begin{lstlisting}[style=ccode, caption={التحقق من مؤشرات المستخدم عند دخول استدعاءات النظام}]
// validate_user_ptr: (*@التحقق من صحة المؤشر قبل الوصول@*)
unsafe fn validate_user_ptr<T>(ptr: *const T) -> \en{Result}<()> {
    if ptr.is_null() {
        return Err(SyscallError::NullPointer);
    }
    let addr = ptr as u64;
    // (*@التحقق من أن العنوان ضمن مجال مساحة المستخدم@*)
    if addr < USER_SPACE_START || addr > USER_SPACE_END {
        return Err(SyscallError::InvalidUserAddress);
    }
    // (*@التحقق من أن الصفحة مُخصَّصة ومعينة للمستخدم@*)
    if !page_table::is_mapped_and_user(addr) {
        return Err(SyscallError::PageNotPresent);
    }
    Ok(())
}
\end{lstlisting}
\end{LTR}

\subsection{خط أنابيب Zero-Copy}

نقل البيانات مباشرة بين العتاد والتطبيق دون نسخ وسيطة يُحقق أداءً استثنائياً:

\begin{center}
\begin{tikzpicture}[
  box/.style={draw=teal, fill=teal!12, rounded corners=3pt,
    minimum width=2.2cm, minimum height=0.8cm, font=\small\bfseries},
  arrow/.style={->, very thick, teal}
]
\node[box] (hw) at (0,0) {Hardware};
\node[box] at (3,0) (dma) {DMA Engine};
\node[box] at (6,0) (shm) {Shared Memory};
\node[box] at (9,0) (app) {Application};

\draw[arrow] (hw) -- (dma) node[midway, above, font=\tiny] {IRQ};
\draw[arrow] (dma) -- (shm) node[midway, above, font=\tiny] {DMA Write};
\draw[arrow] (shm) -- (app) node[midway, above, font=\tiny] {mmap};

\node[font=\small, text=teal, below=0.5cm of shm]
  {\textbf{5.2 GB/s} (شبكة) | \textbf{8.7 GB/s} (تخزين NVMe)};
\end{tikzpicture}
\captionof{figure}{خط أنابيب Zero-Copy: من العتاد إلى التطبيق مباشرة}
\end{center}

\subsection{الدمج Rust/Zig FFI}

سائق \en{blitter} المكتوب بـ Zig يُستدعى من Rust عبر \en{C-ABI}. يمنح هذا النهج:

\begin{itemize}[leftmargin=2em]
\item أداء Zig في مسارات الرسوميات الحرجة (وضع \codeinline{ReleaseFast})
\item ضمانات سلامة Rust في بقية النواة
\item تكامل سلس عبر واجهة \en{C-ABI} الموحدة
\end{itemize}

\subsection{مدير العرض الخطي ومدير النوافذ المتطور (\en{Linear Framebuffer \& NWCC Compositor})}

تم الانتقال من العرض النصي التقليدي إلى عرض رسومي ذي دقة عالية ومحرك نوافذ مستقل مبني بتكامل ثنائي بين Rust وZig:

\begin{itemize}[leftmargin=2em]
\item \textbf{العرض الخطي عالي الدقة (\en{Linear Framebuffer - LFB}):}
تم تمكين وتخصيص مساحة عنوانية خطية لمعالج العرض عالي الدقة، مما أتاح رسم البكسلات مباشرة وتفعيل عمليات الرسم وتغيير المحتوى السريع بدعم من سائق العتاد (\en{hardware blitting}) المطور بلغة Zig لتحقيق كفاءة قصوى للرسوميات ثنائية الأبعاد.

\item \textbf{تحريك النوافذ وجرها تفاعلياً (\en{Window Dragging \& Zig Animation}):}
تم تضمين خوارزميات تفاعلية تسمح للمستخدم بسحب النوافذ (\en{Window Dragging}) وحساب التداخل الرسومي للمناطق بشكل فوري، بالإضافة إلى مكتبة تحريك رسومية (\en{Zig animation}) تعتمد على الفروق الزمنية ووضع المعالجة السريعة لتقديم واجهة مستخدم سلسة وخالية من التقطيع.

\item \textbf{عزل مدير النوافذ عن التطبيقات العميلة (\en{Asynchronous Compositor Task}):}
لتجنب توقف واجهة النوافذ (\en{Compositor Hang}) عند تعليق أو تأخر تطبيق العميل، تم فصل خادم العرض (\en{Compositor}) والعميل (\en{Client}) إلى مهمتين مستقلتين بالكامل (\en{Separate Tasks}) في نظام الجدولة، بحيث يتبادل المعالجان البيانات الرسومية عبر منافذ مخصصة وقنوات اتصال سريعة بشكل مستقل وغير متزامن.
\end{itemize}

\newpage

% ══════════════════════════════════════════════════════════════════
\section{سرعة التطوير والإنتاجية}
% ══════════════════════════════════════════════════════════════════

\subsection{مقاييس الإنتاجية}

تم تطوير النواة خلال 5 أيام فقط بمعدل إنتاجية استثنائية:

\begin{center}
\begin{tabular}{L{5cm}C{3cm}C{3cm}}
\toprule
\rowcolor{titlebg!85}
\bfseries\color{white}\textarabic{المقياس} &
\bfseries\color{white}\textarabic{الإجمالي} &
\bfseries\color{white}\textarabic{المعدل / يوم} \\
\midrule
\rowcolor{rowalt} التزامات (git commits) & 60 & 12 \\
سطور الكود & 28,634 & 5,727 \\
\rowcolor{rowalt} الملفات & 101 & 20 \\
سطور كل التزام & 477 & --- \\
\rowcolor{rowalt} وحدات رئيسية & 8 & 1.6 \\
ميزات مُنجزة & P0+P1 (كاملة) & --- \\
\bottomrule
\end{tabular}
\end{center}

\subsection{جدول زمني للتطوير}

\begin{center}
\begin{tikzpicture}[
  day/.style={draw=accent!60, fill=accent!10, rounded corners=3pt,
    minimum width=2.5cm, minimum height=2.5cm, align=center,
    font=\small, text width=2.2cm},
  arrow/.style={->, thick, accent!60}
]
\node[day, fill=accent!20, draw=accent!70] (1) at (0,0) {\bfseries اليوم 1-2\\[4pt]الهيكل الأساسي\\ الذاكرة\\ x86\_64 معمارية};
\node[day, fill=accent2!20, draw=accent2!70] (2) at (3.0,0) {\bfseries اليوم 2-3\\[4pt]نظام القدرات\\ جدولة MLFQ\\ طبقة \en{ABI}};
\node[day, fill=teal!20, draw=teal!70] (3) at (6.0,0) {\bfseries اليوم 3-4\\[4pt]eBPF VM والمدقق\\ نظام ZiqaFS\\ VFS};
\node[day, fill=orange!20, draw=orange!70] (4) at (9.0,0) {\bfseries اليوم 4-5\\[4pt]VirtIO, VGA, PCI\\ Shell (40+ أمر)\\ Tetris \& DOOM};
\foreach \s/\t in {1/2, 2/3, 3/4}{
  \draw[arrow] (\s.east) -- (\t.west);
}
\end{tikzpicture}
\captionof{figure}{الجدول الزمني لتطوير \textsc{ZiqaKernel} عبر 5 أيام}
\end{center}

\vspace{0.3cm}

\begin{center}
\begin{tikzpicture}
\begin{axis}[
  x dir=reverse,
  width=0.85\textwidth, height=6cm,
  xlabel={اليوم},
  ylabel={سطور الكود التراكمية},
  xmin=0, xmax=6, ymin=0, ymax=32000,
  xtick={1,2,3,4,5},
  ymajorgrids=true,
  grid style={dashed, gray!25},
  smooth,
]
\addplot[color=accent, thick, mark=*,
  mark options={fill=accent}] coordinates {
  (0,0) (1,5200) (2,11400) (3,17800) (4,23500) (5,28633)
};
\addplot[color=accent3, dashed, thick] coordinates {
  (0,0) (5,28633)
};
\node[font=\small, text=accent] at (axis cs:4.5,26000) {نمو فعلي};
\node[font=\small, text=accent3] at (axis cs:4,19000) {نمو خطي};
\end{axis}
\end{tikzpicture}
\captionof{figure}{نمو سطور الكود التراكمية عبر 5 أيام من التطوير}
\end{center}

\newpage

% ══════════════════════════════════════════════════════════════════
\section{طبقة استدعاءات النظام (\en{ABI} Layer)}
% ══════════════════════════════════════════════════════════════════

تدعم النواة 129 استدعاء نظام موزعة عبر 4 طبقات متسلسلة:

\begin{center}
\begin{tikzpicture}[
  layer/.style={draw, rounded corners=5pt, minimum width=10cm,
    minimum height=0.9cm, font=\bfseries\small, align=center},
  count/.style={font=\small, text=darkgray}
]
\node[layer, fill=accent2!20, draw=accent2] at (0,3)
  {Linux \en{ABI} \en{Plugin} --- 99 استدعاء (\en{POSIX} متوافق)};
\node[count] at (5.8,3) {متوافق Linux};

\node[layer, fill=accent!20, draw=accent] at (0,2)
  {Core Native --- 25 (قدرات، IPC، \en{KVM})};
\node[count] at (5.8,2) {أصيل};

\node[layer, fill=orange!20, draw=orange] at (0,1)
  {Fast-path --- 4 (shm\_get, shm\_at, io\_uring\_*)};
\node[count] at (5.8,1) {مسار سريع};

\node[layer, fill=purple!15, draw=purple] at (0,0)
  {\en{WASM} \en{ABI} --- 1 (proc\_exit)};
\node[count] at (5.8,0) {WASM};

% اسم المجموعة الكلية
\node[font=\large\bfseries, text=titlebg] at (0,-0.8)
  {إجمالي: 129 استدعاء نظام};

\end{tikzpicture}
\captionof{figure}{تسلسل طبقات استدعاءات النظام في \textsc{ZiqaKernel}}
\end{center}

\vspace{0.3cm}

ينفذ كل استدعاء عبر \codeinline{int 0x80} مع جدول توزيع قابل للتوسيع:

\begin{equation}
\text{بنية السياق} = \Big(\underbrace{\mathrm{rax}}_{\text{رقم}},\;
\underbrace{\mathrm{rdi}}_{\text{arg}_1},\;
\underbrace{\mathrm{rsi}}_{\text{arg}_2},\;
\underbrace{\mathrm{rdx}}_{\text{arg}_3},\;
\underbrace{\mathrm{r10}}_{\text{arg}_4},\;
\underbrace{\mathrm{r8}}_{\text{arg}_5},\;
\underbrace{\mathrm{r9}}_{\text{arg}_6}\Big) + \text{PID Space}
\end{equation}

\subsection{توسيع استدعاءات التوافقية (\en{POSIX \& IPC Expansion})}

لتسهيل تشغيل برمجيات مساحة المستخدم القياسية (مثل \en{BusyBox})، تم توسيع استدعاءات نظام التوافقية ومنافذ التواصل الأصيل:

\begin{itemize}[leftmargin=2em]
\item \textbf{استدعاءات التحكم في العمليات والملفات (\en{POSIX Syscalls}):}
  \begin{itemize}
  \item \codeinline{sys\_waitpid}: تمت ترقية الاستدعاء ليقوم بنقل حالة الخروج للعملية الابن بشكل متزامن وحقنها في مرجع الحالة الممرر من قبل عملية الأب، مع ضبط خيارات الانتظار المتزامنة.
  \item \codeinline{sys\_stat}: يوفر استعلاماً كاملاً لخصائص الملفات وحجمها وأذونات الوصول إليها مباشرة من نظام VFS.
  \item \codeinline{sys\_pipe}: ينشئ قناة اتصال ثنائية الاتجاه (\en{IPC Pipe}) لنقل البيانات بين العمليات بشكل متسلسل ودون الحاجة لتخزين وسيط على القرص.
  \end{itemize}
\item \textbf{استدعاءات الذاكرة المشتركة (\en{Shared Memory System}):}
تم توحيد وزيادة استدعاءات الذاكرة المشتركة للأداء العالي (SHM) مثل \codeinline{sys\_shmget} و\codeinline{sys\_shmat} لدعم النقل المباشر للمقاطع الكبيرة بين العمليات بدون نسخ، ومزامنة واجهة تطبيقات Zig لضمان تطابق سياق فضاء المستخدم مع سياق النواة.
\end{itemize}

\newpage

% ══════════════════════════════════════════════════════════════════
\section{الخاتمة والعمل المستقبلي}\label{sec:conc}
% ══════════════════════════════════════════════════════════════════

\subsection{الخلاصة}

نواة \textsc{ZiqaKernel} تُثبت أن بناء نظام تشغيل حديث آمن وفعّال ومرن أمر ممكن بلغة \textsc{Rust}. خلال 5 أيام من البحث المكثف، تحققت جميع الأهداف المخططة:

\begin{resultbox}[ملخص الإنجازات]
\begin{itemize}[leftmargin=2em]
\item \textbf{نظام أمان \en{A+ Tier}:} إلغاء فوري ومتكرر، شجرة تفويض متعددة المستويات، عزل \en{Ring 3} كامل مع SMEP/SMAP/UMIP المتكامل.
\item \textbf{جدولة وتحصين \en{A-Tier}:} جدولة \en{MLFQ} مع 4 طوابير أولوية ديناميكية، مع تحصين كامل للمجدول ضد المقاطعات المتزامنة (\en{Interrupt-Safe scheduling}).
\item \textbf{\en{eBPF Obsidian-Tier}:} مدقق تحليل تجريدي، تتبع أنواع كامل، حلقات مقيدة (\en{bounded loops})، وخرائط التجزئة والبرمجيات (\en{Hash \& ProgArray})، مع 11 مساعداً متطوراً ومكدس آلة افتراضية مستقل.
\item \textbf{\en{ABI} كامل بالتوافقية:} طبقة استدعاءات نظام موحدة تدعم 129 استدعاء نظام، منها استدعاءات متكاملة للتوافق مع Linux لتشغيل BusyBox بالإضافة إلى استدعاءات الذاكرة المشتركة الأصيلة.
\item \textbf{تطبيقات وتجربة بصرية:} خادم عرض نوافذ مستقل (\en{NWCC Compositor})، لوحة عرض خطية عالية الدقة (\en{LFB}) مع محرك رسم بلغة Zig تفاعلي، وتطبيقات رسومية (\en{Tetris} و\en{DOOM}) على المعدن العاري.
\item \textbf{خريطة معرفة:} 2,015 عقدة، 3,843 علاقة، 131 مجتمعاً معمارياً.
\end{itemize}
\end{resultbox}

\subsection{الأولويات المستقبلية}

\begin{center}
\begin{tabular}{C{1.2cm}C{2.5cm}L{8cm}}
\toprule
\rowcolor{titlebg!85}
\bfseries\color{white}\textarabic{الأولوية} &
\bfseries\color{white}\textarabic{المجال} &
\bfseries\color{white}\textarabic{الوصف} \\
\midrule
\rowcolor{rowalt!60} \en{P2} & \en{ARM64/RISC-V} &
  دعم معماريات متعددة: إعادة كتابة GDT/IDT لـ GIC، ASM Trampoline، MMU \\
\en{P2} & \en{Cloud} & تكامل سحابي، \en{KVM} \en{hypervisor} كامل، صور قابلة للنشر \\
\rowcolor{rowalt!60} \en{P2} & \en{Plugin} \en{ABI} &
  توسيع طبقة \en{ABI} لدعم \en{WASM} كامل، eBPF \en{JIT compiler} \\
\en{P2} & \en{Dev Tools} &
  مُصحِّح (\en{debugger}) مُدمَج، تتبع كامل، أدوات تشخيص أداء \\
\rowcolor{rowalt!60} \en{P3} & \en{Network} &
  \en{TCP/IP} stack كامل، دعم \en{TLS}، بروتوكولات إنترنت الأشياء \\
\en{P3} & \en{SMP Tuning} &
  تحسين الأداء لـ 32+ نواة، \en{NUMA awareness}، توازن حمل ديناميكي \\
\rowcolor{rowalt!60} \en{P3} & التحقق الرسمي &
  توظيف \en{Isabelle/HOL} لتحقق رسمي من نظام القدرات وMFLQ \\
\bottomrule
\end{tabular}
\end{center}

\vspace{0.5cm}

\begin{warningbox}[القيود الحالية]
النواة لا تزال تجريبية وغير مناسبة للإنتاج. القيود الرئيسية:
\begin{enumerate}[leftmargin=2em]
\item دعم معمارية واحدة فقط (x86\_64) حالياً
\item نظام eBPF مبسط مقارنة بنظام \en{Linux} (رغم توسيعه ودعمه لخرائط التجزئة والحلقات المقيدة)
\item غياب التحقق الرسمي الكامل (على خلاف seL4)
\item نظام الملفات \en{ZiqaFS} أساسي (بدون \en{journaling} كامل)
\end{enumerate}
\end{warningbox}

\newpage

% ══════════════════════════════════════════════════════════════════
\section*{المراجع}
% ══════════════════════════════════════════════════════════════════

\begin{thebibliography}{10}

\bibitem{graphify}
\en{S}.~Shamsi,
``\en{graphify}: Knowledge Graph Pipeline for Code and Document Corpora,''
2026. \url{https://github.com/safishamsi/graphify}

\bibitem{ebpf}
M.~Fleming,
``A Thorough Introduction to eBPF,''
\textit{\en{Linux} Weekly News}, 2020.

\bibitem{rust}
\en{S}.~Klabnik and C.~Nichols,
\textit{The Rust Programming Language}, 2nd~ed.,
No Starch Press, 2022.

\bibitem{zig}
A.~Kelley,
``Zig Programming Language,''
\textit{ziglang.org}, 2023.

\bibitem{seL4}
G.~Klein et al.,
``seL4: Formal Verification of an OS Kernel,''
\textit{Proc. SOSP}, 2009.

\bibitem{mlfq}
J.~K.~Ousterhout,
``Scheduling Techniques for Concurrent Systems,''
\textit{Proc. ICDCS}, 1982.

\bibitem{capability}
J.~\en{S}.~Shapiro,
``EROS: A \en{Capability-Based} Operating System,''
Ph.D. dissertation, University of Pennsylvania, 1999.

\bibitem{amdahl}
G.~M.~Amdahl,
``Validity of the Single Processor Approach to Achieving Large-Scale
Computing Capabilities,''
\textit{Proc. AFIPS Spring Joint Computer Conference}, 1967.

\bibitem{abstract}
P.~Cousot and R.~Cousot,
``\en{Abstract Interpretation}: A Unified Lattice Model for Static Analysis
of Programs,''
\textit{Proc. POPL}, 1977.

\bibitem{louvain}
V.~D.~Blondel et al.,
``F\en{ast} unfolding of communities in large networks,''
\textit{Journal of Statistical Mechanics}, 2008.

\end{thebibliography}

\vspace{0.8cm}
\begin{center}
\rule{0.5\textwidth}{0.4pt}\\[6pt]
{\small\color{darkgray}
تم الإنشاء باستخدام
\textbf{graphify} knowledge graph pipeline \;|\;
\LaTeX{} (XeLaTeX) \;|\;
FreeSerif \;|\;
pgfplots + tikz-3dplot
}\\[4pt]
{\small\color{darkgray}
البيانات: \codeinline{\en{graphify}-out/graph.json} \;|\;
المصدر: \codeinline{github.com/\en{ZiqaKernel}/my-os}
}
\end{center}

\end{document}